\chapter{Statistics}
\label{ch:statistics}

\begin{tcolorbox}[feynbox={Sky}{BBg},title={Before and After}]
\textbf{Before this chapter you need:} logarithms and the log rules
(\S\ref{sec:logrules}), derivatives and how to find a maximum
(\S\ref{sec:rules}--\S\ref{sec:secondderiv}), gradients (\S\ref{sec:gradient}),
random variables, expectation and variance
(\S\ref{sec:randomvars}--\S\ref{sec:variance}), the law of large numbers and the
CLT (\S\ref{sec:lln}--\S\ref{sec:clt}).\\
\textbf{After it you will understand:} why an estimate is itself a random
quantity, what maximum likelihood maximises, why this project fits GARCH with a
Gaussian likelihood it knows to be wrong, what a $p$-value does and does not
mean, why a test on 125 days is nearly powerless, and what goes wrong when 19
metrics are tested at once.
\end{tcolorbox}

Probability reasons from a known distribution to the behaviour of data.
Statistics runs the arrow backwards: from data in hand to claims about the
distribution that produced it. Every number reported in this thesis is a
statistic, and every one carries uncertainty this chapter makes explicit.

%% ─────────────────────────────────────────────────────────────────────────────
\section{Samples and Sampling Distributions}
\label{sec:sampling}

\situation{
You measure the temperature every day for a week: 15, 17, 14, 19, 16, 21, 13
degrees. You want the ``typical'' value and how much the values vary. Next week
you repeat the exercise and get different numbers --- and therefore a different
average.
}

\problem{
The second week's average differs from the first, yet neither is wrong. The
summary you compute is not a fixed property of the world but something that
changes with the sample. Unless that variation is quantified, there is no way to
tell a real difference from an accident of sampling.
}

\workedarith{
\textbf{The week's data.} $x = \{15, 17, 14, 19, 16, 21, 13\}$.

\medskip
\textbf{Mean:} $\bar{x} = (15+17+14+19+16+21+13)/7 = 115/7 = 16.43$.

\textbf{Deviations:} $\{-1.43, 0.57, -2.43, 2.57, -0.43, 4.57, -3.43\}$.

\textbf{Squared:} $\{2.04, 0.33, 5.90, 6.60, 0.18, 20.88, 11.76\}$, totalling $47.69$.

\textbf{Sample variance:} $s^2 = 47.69/6 = 7.95$, so $s = 2.82$\textdegree C.

\medskip
\textbf{Why divide by $6$ and not $7$.} The deviations are taken from
$\bar{x}$, which was computed from the same seven numbers. They therefore sum to
zero by construction, so only six of them are free --- fix any six and the seventh
follows. Dividing by $n-1$ corrects for that, and without it the variance would be
systematically too small.

\medskip
\textbf{The sampling distribution.} If temperatures have true standard deviation
$\sigma = 3$, then across repeated weeks the sample mean has standard deviation
\[
  \frac{\sigma}{\sqrt{n}} = \frac{3}{\sqrt{7}} = 1.13 .
\]
So a second week averaging $17.5$ instead of $16.4$ is entirely unremarkable: the
gap of $1.1$ is one standard error.
}

\formaldef{
A \textbf{sample} is $n$ observations drawn from a population. A
\textbf{statistic} is any function of the sample; being a function of random
quantities, it is itself a random variable, and its distribution across repeated
samples is its \textbf{sampling distribution}.

\[
  \bar{x} = \frac{1}{n}\sum_{i=1}^n x_i , \qquad
  s^2 = \frac{1}{n-1}\sum_{i=1}^n (x_i - \bar{x})^2 .
\]
The \textbf{standard error} of the mean is $\sigma/\sqrt{n}$, estimated by
$s/\sqrt{n}$. The $n-1$ divisor is \textbf{Bessel's correction}.
}

\analogybreak{
Standard deviation summarises spread well only for roughly symmetric
distributions. Returns fall much further than they rise, so a single number
understates downside risk. And the $\sigma/\sqrt{n}$ standard error is itself a
CLT result requiring finite variance --- available for returns
($\hat\alpha > 2$) but not for their fourth moment, which is why no standard error
is ever quoted for a sample kurtosis in this thesis.
}

\whyhere{
Daily volatility $\sigma = \mathrm{sd}(r_t)$ is the foundational scale in all five
markets: MOEX $1.87\%$/day, BOVESPA $1.63\%$. Annualising multiplies by
$\sqrt{252}$. Every reported metric is a statistic with a sampling distribution,
which is why the pooled test set of about $2{,}480$ observations is preferred to
the per-market $496$ wherever the metric permits pooling.
}

%% ─────────────────────────────────────────────────────────────────────────────
\section{Estimators: Bias and Variance}
\label{sec:estimators}

\situation{
Two bathroom scales. One reads exactly $2$\,kg heavy every time --- consistent but
wrong. The other averages correctly across many weighings but scatters
$\pm 3$\,kg. Neither is obviously better; it depends what you need.
}

\problem{
``Accurate'' bundles together two distinct failures --- being systematically off,
and being erratic --- and they trade against each other. Separating them is what
lets you choose an estimator deliberately rather than by habit.
}

\workedarith{
True weight $70$\,kg. Five weighings on each scale.

\medskip
\begin{tabular}{@{}lrrrrrrr@{}}
\toprule
& \multicolumn{5}{c}{readings} & mean & spread \\
\midrule
Scale A & $72$ & $72$ & $72$ & $72$ & $72$ & $72$ & $0$ \\
Scale B & $67$ & $73$ & $70$ & $68$ & $72$ & $70$ & $2.55$ \\
\bottomrule
\end{tabular}

\medskip
\textbf{Scale A:} bias $= 72 - 70 = +2$, variance $= 0$.

\textbf{Scale B:} bias $= 0$, standard deviation $= 2.55$.

\medskip
\textbf{Which is better?} Compare mean squared error, $\mathrm{MSE} =
\text{bias}^2 + \text{variance}$:
\[
  \text{A: } 2^2 + 0 = 4 , \qquad \text{B: } 0^2 + 6.5 = 6.5 .
\]
The biased scale wins --- on this criterion, at this sample size. Averaging four
readings from B, however, cuts its variance to $6.5/4 = 1.63$ and its MSE to
$1.63$, beating A. The right choice depends on how much data you have.

\medskip
\textbf{Bessel's correction as bias removal.} On the seven temperatures,
dividing by $n = 7$ gives $47.69/7 = 6.81$; by $n - 1 = 6$ gives $7.95$. The first
is biased low; the second is unbiased.
}

\formaldef{
An \textbf{estimator} $\hat\theta$ is a statistic used to estimate a parameter
$\theta$. Its \textbf{bias} is $\E[\hat\theta] - \theta$, and it is
\textbf{unbiased} when this is zero.

\[
  \mathrm{MSE}(\hat\theta) = \E\bigl[(\hat\theta - \theta)^2\bigr]
  = \bigl(\text{bias}\bigr)^2 + \Var(\hat\theta) .
\]

$\hat\theta$ is \textbf{consistent} if it converges to $\theta$ as $n \to \infty$.
Consistency is about the limit; unbiasedness is about any fixed $n$. Neither
implies the other.
}

\analogybreak{
Unbiasedness is a weaker virtue than its name suggests. It guarantees correctness
on average over hypothetical repetitions, which is cold comfort when you have one
sample. A biased estimator with much smaller variance is frequently the better
instrument, and every regularised model in machine learning --- including weight
decay and early stopping --- deliberately accepts bias to reduce variance.
}

\whyhere{
The Hill estimator is biased at small samples yet consistent, and its bias shrinks
as the tail sample grows. This is why Chapter~\ref{ch:heavytails} reports it
converging from $4.83$ at $n=250$ to $3.11$ at $n=2{,}000$ on BOVESPA: the
movement is bias disappearing, not instability. Sample kurtosis over the same
range moves the opposite way --- $2.53$ to $11.84$ --- because it is not converging
to anything at all.
}

%% ─────────────────────────────────────────────────────────────────────────────
\section{Maximum Likelihood}
\label{sec:mle}

\situation{
A coin lands heads 7 times in 10 flips. What is its bias? You would guess $0.7$,
and the reasoning behind that guess is: of all possible biases, $0.7$ is the one
that makes what actually happened most probable.
}

\problem{
That reasoning needs to become a procedure that works for any model, not just
coins --- one that turns ``which parameter best explains this data'' into a
calculation.
}

\workedarith{
\textbf{The coin.} For bias $p$, the probability of a specific sequence with 7
heads in 10 flips is $p^7(1-p)^3$. Evaluate:

\medskip
\begin{tabular}{@{}lrrrrr@{}}
\toprule
$p$ & $0.5$ & $0.6$ & $0.7$ & $0.8$ & $0.9$ \\
\midrule
$p^7(1-p)^3$ & $0.000977$ & $0.001791$ & $\mathbf{0.002224}$ & $0.001678$ & $0.000478$ \\
\bottomrule
\end{tabular}

\medskip
The largest value is at $p = 0.7$. \checkmark

\medskip
\textbf{By calculus.} Maximise the logarithm instead --- the log is increasing, so
it has its maximum in the same place, and it turns the product into a sum
(\S\ref{sec:logrules}):
\[
  \ell(p) = 7\ln p + 3\ln(1-p) .
\]
Differentiate and set to zero, using $\frac{d}{dp}\ln p = 1/p$:
\[
  \ell'(p) = \frac{7}{p} - \frac{3}{1-p} = 0
  \;\Longrightarrow\; 7(1-p) = 3p
  \;\Longrightarrow\; 7 = 10p
  \;\Longrightarrow\; p = 0.7 . \quad\checkmark
\]

\textbf{Confirm it is a maximum.} $\ell''(p) = -7/p^2 - 3/(1-p)^2 < 0$ for every
$p$ in $(0,1)$, so the function is concave and $0.7$ is its highest point
(\S\ref{sec:secondderiv}). \checkmark
}

\formaldef{
Given data $x_1,\ldots,x_n$ and a model with density $f(x \mid \theta)$, the
\textbf{likelihood} treats the data as fixed and $\theta$ as the argument:
\[
  L(\theta) = \prod_{i=1}^n f(x_i \mid \theta) ,
  \qquad
  \ell(\theta) = \ln L(\theta) = \sum_{i=1}^n \ln f(x_i \mid \theta) .
\]
The \textbf{maximum likelihood estimator} is $\hat\theta = \arg\max_\theta \ell(\theta)$.

Under regularity conditions --- including that the model is correctly specified ---
MLE is consistent, asymptotically unbiased, and asymptotically normal with the
smallest achievable variance.
}

\analogybreak{
The likelihood is not a probability distribution over $\theta$. It is a function
of $\theta$ computed from fixed data, and it does not integrate to one; reading
$L(0.7) = 0.0022$ as ``the probability the bias is $0.7$'' is a category error.
Turning the likelihood into a distribution over $\theta$ requires a prior and
Bayes' theorem (\S\ref{sec:bayes}).

The optimality guarantees also assume the model is right. When it is not, the
next section describes what survives.
}

\whyhere{
Log-likelihood is maximised whenever a GARCH model is fitted in this project. The
sum form is not a convenience but a necessity: a product of several thousand small
densities underflows to zero in floating-point arithmetic, while the corresponding
sum of logarithms is perfectly stable.
}

%% ─────────────────────────────────────────────────────────────────────────────
\section{QMLE: Consistency Under Misspecification}
\label{sec:qmle}

\situation{
You measure a room with a tape that stretches slightly under tension. Every
individual reading is a little off. But if the stretch is proportional, the
\emph{ratio} of two walls is still exactly right --- the flawed instrument gets
some quantities wrong while getting others right.
}

\problem{
Maximum likelihood's guarantees assume the model is correct. Real models never
are. The practical question is not whether the model is wrong --- it is --- but
which of its outputs survive being wrong, and this project's residual-kurtosis
test depends entirely on the answer.
}

\workedarith{
\textbf{The setup.} GARCH$(1,1)$ models the conditional variance:
\[
  \sigma_t^2 = \omega + \alpha \varepsilon_{t-1}^2 + \beta \sigma_{t-1}^2 ,
  \qquad \varepsilon_t = \sigma_t z_t ,
\]
where $z_t$ is a standardised innovation. Fitting requires assuming a
distribution for $z_t$.

\medskip
\textbf{The test we want to run.} Real returns keep excess kurtosis \emph{after}
dividing out the time-varying volatility (Bollerslev 1987). That is stylised
fact 7: fat tails are not fully explained by volatility clustering. We want to
measure that leftover kurtosis in the standardised residuals
$\hat z_t = \varepsilon_t / \hat\sigma_t$ and check whether a generator
reproduces it.

\medskip
\textbf{Why not assume a Student-$t$?} A $t$ innovation would fit the data better
--- the returns really are heavy-tailed. But it fits better precisely by absorbing
the excess kurtosis into the assumed innovation distribution. The residuals would
then look $t$-shaped by construction, and measuring their kurtosis would recover
the assumption rather than a property of the data. The test would be circular.

\medskip
\textbf{Why Gaussian works anyway.} Maximising a Gaussian likelihood estimates
$(\omega, \alpha, \beta)$ consistently even when $z_t$ is not Gaussian. The
Gaussian score equations depend on the data only through first and second
conditional moments, and those are correctly specified by the GARCH recursion
whatever the shape of $z_t$. So:
\begin{itemize}
  \item the variance path $\hat\sigma_t$ is estimated consistently --- \emph{right};
  \item the standard errors from the Gaussian information matrix are --- \emph{wrong},
        and need a sandwich correction;
  \item leftover kurtosis in $\hat z_t$ is --- \emph{genuine evidence}, because
        nothing in the fitting procedure put it there.
\end{itemize}

The deliberately wrong assumption is what makes the measurement honest.
}

\formaldef{
\textbf{Quasi-maximum likelihood estimation (QMLE)} maximises a likelihood known
to be misspecified. For GARCH, White (1982) and Bollerslev--Wooldridge (1992)
establish that the Gaussian QMLE of the conditional-variance parameters is
consistent and asymptotically normal under mild conditions, whatever the true
innovation distribution.

The asymptotic variance is \emph{not} the usual inverse information matrix but the
\textbf{sandwich} form
\[
  A^{-1} B A^{-1} ,
\]
where $A$ is the expected Hessian and $B$ the variance of the score. Under correct
specification $A = B$ and this collapses to the standard result.
}

\analogybreak{
Consistency under misspecification is a guarantee about specific parameters, not a
licence to ignore the model being wrong. Gaussian QMLE gets the variance dynamics
right and the uncertainty quantification wrong. It also assumes the \emph{variance
equation} is correctly specified --- if the true process has leverage effects or
long-memory volatility that GARCH$(1,1)$ cannot represent, no choice of innovation
distribution repairs that.
}

\whyhere{
This is the reasoning behind \texttt{resid\_kurtosis\_diff}, one of the seven
ranked temporal metrics. The pipeline fits a constant-mean GARCH$(1,1)$ by
Gaussian QMLE, extracts $\hat z_t = r_t/\hat\sigma_t$, and compares the excess
kurtosis of real and synthetic residuals. A generator that reproduces volatility
clustering but not conditional fat tails will match the ACF metrics and fail this
one --- which is exactly the discrimination the metric exists to provide.
}

%% ─────────────────────────────────────────────────────────────────────────────
\section{Hypothesis Testing}
\label{sec:hypothesis}

\situation{
A colleague claims they can tell tap water from bottled. You pour ten pairs and
they get eight right. Impressive --- or is it? Someone guessing would average five,
and getting eight by luck alone is not extraordinary.
}

\problem{
Distinguishing a real effect from a lucky run needs a rule fixed in advance.
Without one, any result can be narrated as impressive after the fact.
}

\workedarith{
\textbf{Null hypothesis:} they are guessing, so each trial is a fair coin,
$p = 0.5$.

Probability of 8 or more correct out of 10 by chance:
\[
  \Prob(8) = \binom{10}{8}(0.5)^{10} = 45/1024 = 0.0439
\]
\[
  \Prob(9) = \binom{10}{9}(0.5)^{10} = 10/1024 = 0.0098 , \qquad
  \Prob(10) = 1/1024 = 0.0010
\]
\[
  \Prob(\geq 8) = (45 + 10 + 1)/1024 = 56/1024 = 0.0547 .
\]

\medskip
At the conventional $5\%$ threshold, $0.0547 > 0.05$: do not reject. Eight out of
ten is not quite enough.

\textbf{Nine out of ten:} $\Prob(\geq 9) = 11/1024 = 0.0107 < 0.05$. Reject.

\medskip
The line between "unconvincing" and "convincing" falls between 8 and 9 correct ---
a single trial. That sharpness is an artefact of the threshold, not a feature of
the evidence.
}

\formaldef{
A test contrasts a \textbf{null hypothesis} $H_0$ with an \textbf{alternative}
$H_1$. A \textbf{test statistic} is computed, and its distribution under $H_0$
determines whether the observed value is surprising.

Two errors are possible:
\begin{itemize}
  \item \textbf{Type I}: rejecting a true $H_0$. Probability $\alpha$, chosen in
        advance, conventionally $0.05$.
  \item \textbf{Type II}: failing to reject a false $H_0$. Probability $\beta$;
        the \textbf{power} is $1 - \beta$ (\S\ref{sec:power}).
\end{itemize}
Failing to reject is not evidence for $H_0$. It means the data was insufficient to
rule it out.
}

\analogybreak{
The $0.05$ threshold is a convention with no mathematical standing --- Fisher
proposed it as a rough guide, not a law. Treating $0.049$ and $0.051$ as
categorically different, when they differ by less than the sampling noise in
either, is the single most common abuse of the framework.
}

\whyhere{
The Kupiec test asks whether an observed VaR exceedance rate is distinguishable
from the claimed $5\%$; Christoffersen asks whether exceedances cluster; the
Anderson--Darling and Kolmogorov--Smirnov tests compare distributions. Each has a
null, a statistic and a threshold --- and \S\ref{sec:power} shows that on
per-market samples several of them cannot reject anything.
}

%% ─────────────────────────────────────────────────────────────────────────────
\section{What a $p$-Value Is and Is Not}
\label{sec:pvalues}

\situation{
A study reports $p = 0.03$ and the write-up says ``there is a 3\% chance the
result is due to chance''. That sentence is wrong, and the correct statement is
narrower and less satisfying than most readers want.
}

\problem{
The $p$-value is the most misread number in quantitative work. Since several
appear in this thesis --- Kupiec, Christoffersen, ARCH, Anderson--Darling --- what
each one licenses needs stating precisely.
}

\workedarith{
Return to eight correct out of ten. The $p$-value is $0.0547$.

\medskip
\textbf{What it means:} \emph{if} the person were guessing, the probability of
seeing eight or more correct is $0.0547$.

\medskip
\textbf{What it does not mean:}
\begin{itemize}
  \item Not the probability they are guessing. That is
        $\Prob(H_0 \mid \text{data})$, and getting it from
        $\Prob(\text{data} \mid H_0)$ requires a prior and Bayes' theorem
        (\S\ref{sec:bayes}) --- the two differ by exactly the factor that made a
        positive disease test only $9\%$ informative.
  \item Not the probability the result is a fluke.
  \item Not a measure of effect size. With $n = 10{,}000$ a trivial edge such as
        $50.5\%$ produces a tiny $p$; with $n = 10$ a large edge produces a big one.
  \item Not comparable across tests as a strength ranking.
\end{itemize}

\medskip
\textbf{A concrete demonstration of the last point.} Someone correct on $51\%$ of
$100{,}000$ trials gives $p \approx 0.00003$ --- overwhelming significance for an
edge nobody would care about. Someone correct on $90\%$ of $10$ trials gives
$p = 0.0107$. The second is the far more impressive ability and the less
significant result.
}

\formaldef{
The \textbf{$p$-value} is the probability, computed under $H_0$, of observing a
test statistic at least as extreme as the one obtained:
\[
  p = \Prob\bigl(T \geq t_{\text{obs}} \mid H_0\bigr) .
\]
It is a statement about data given a hypothesis, never about a hypothesis given
data. Under $H_0$ the $p$-value is uniformly distributed on $[0,1]$ --- which is
why one test in twenty falls below $0.05$ when nothing whatever is going on.
}

\analogybreak{
Some $p$-values in this thesis are not even usable as continuous evidence.
\texttt{scipy}'s \texttt{anderson\_ksamp} clamps its result to the range
$[0.001, 0.25]$ and warns whenever the true value falls outside. A reported
$0.001$ may mean $0.001$ or $10^{-12}$; the number has been censored. The
\texttt{arch\_pvalue\_diff} metric was moved to descriptive-only partly for this
reason --- it saturates on real data, giving $|\Delta p| = 0.000000$ and no
discrimination at all.
}

\whyhere{
Kupiec and Christoffersen $p$-values are reported per model. They are read as
``the data does not contradict correct coverage'', never as ``the model is
correct'' --- and given the power calculations in the next section, on per-market
samples the first statement is often true simply because almost nothing could
contradict it.
}

%% ─────────────────────────────────────────────────────────────────────────────
\section{Confidence Intervals}
\label{sec:confint}

\situation{
A poll reports 52\% support, ``margin of error 3 points''. The single number 52 is
less informative than the range 49--55, because the range carries its own
uncertainty with it.
}

\problem{
A point estimate invites false precision. Reporting $\hat\alpha = 3.11$ suggests
a confidence that a sample of a few hundred tail observations cannot support, and
an interval says so honestly.
}

\workedarith{
The seven temperatures: $\bar{x} = 16.43$, $s = 2.82$, $n = 7$.

\medskip
\textbf{Standard error:} $s/\sqrt{n} = 2.82/\sqrt{7} = 1.066$.

\textbf{95\% interval} (using $\approx 2$ standard errors):
\[
  16.43 \pm 2 \times 1.066 = 16.43 \pm 2.13 = [14.30,\; 18.56] .
\]

\medskip
\textbf{The effect of more data.} With the same $s$ but $n = 700$:
\[
  \text{SE} = 2.82/\sqrt{700} = 0.107 , \qquad
  \text{interval} = 16.43 \pm 0.21 .
\]
A hundredfold increase in data narrows the interval tenfold --- the
$\sqrt{n}$ of \S\ref{sec:lln} again.

\medskip
\textbf{The discriminative AUC null.} Fifteen real-versus-real splits gave
$0.506 \pm 0.084$. An observed AUC of $0.62$ sits
\[
  z = \frac{0.62 - 0.506}{0.084} = 1.36
\]
standard errors above the null centre, giving $p \approx 0.17$. Not significant
--- despite $0.62$ looking comfortably above $0.5$ to the naked eye.
}

\formaldef{
A \textbf{$(1-\alpha)$ confidence interval} is constructed so that, across
repeated samples, the interval contains the true parameter a proportion $1-\alpha$
of the time. For a mean with known $\sigma$:
\[
  \bar{x} \pm z_{\alpha/2}\,\frac{\sigma}{\sqrt{n}} ,
  \qquad z_{0.025} = 1.96 .
\]
With $\sigma$ estimated by $s$ and small $n$, the $t$-distribution replaces the
normal.
}

\analogybreak{
The confidence is in the procedure, not in the particular interval. Once computed,
$[14.30, 18.56]$ either contains the true mean or it does not --- there is no
probability left. ``95\% probability the mean lies in this interval'' is a
Bayesian credible interval, which requires a prior and is a different object.

Note also that the interval above assumes the CLT applies. For a statistic with no
finite variance --- sample kurtosis on our data --- the construction is unavailable
in principle, not merely difficult.
}

\whyhere{
The $0.506 \pm 0.084$ null is why the evaluation reports
\texttt{discriminative\_auc\_absz}, the distance from the null centre measured in
standard errors, rather than the raw AUC. Judging a generator against a
theoretical $0.5$ would flag ordinary sampling noise as a detectable difference.
}

%% ─────────────────────────────────────────────────────────────────────────────
\section{Statistical Power}
\label{sec:power}

\situation{
You test a bathroom scale for a 200\,g bias by weighing yourself twice. The test
will not detect it --- not because the bias is absent, but because two weighings
cannot resolve 200\,g against day-to-day variation. The instrument is too blunt
for the question.
}

\problem{
A test that fails to reject tells you one of two things: the null is true, or your
test was too weak to notice. These are entirely different conclusions, and only a
power calculation separates them.
}

\workedarith{
\textbf{Kupiec, this project's numbers.} A VaR model claims $5\%$ exceedances.
Suppose the truth is $7\%$. Can the test detect the difference?

\medskip
\begin{tabular}{@{}lrl@{}}
\toprule
Test days $n$ & Power at true $p = 7\%$ & Setting \\
\midrule
$125$    & $17\%$ & 5-year, per market \\
$496$    & $\approx 56\%$ & 20-year, per market \\
$2{,}480$ & $\approx 99\%$ & 20-year, pooled across 5 markets \\
\bottomrule
\end{tabular}

\medskip
\textbf{Reading the top row.} With 125 days, a genuinely broken model --- exceeding
its VaR $40\%$ more often than advertised --- goes undetected five times in six. A
non-rejection carries essentially no information.

\medskip
\textbf{Why pooling helps so much.} Exceedance counts scale with $n$ while their
standard deviation scales with $\sqrt{n}$, so the signal-to-noise ratio grows as
$\sqrt{n}$. Going from $496$ to $2{,}480$ multiplies it by $\sqrt{5} = 2.24$, and
that is enough to move power from a coin flip to near certainty.

\medskip
\textbf{The asymmetry this creates.} At $n = 125$ a non-rejection is
uninformative, while a rejection still means something --- Type I error is
controlled at $5\%$ regardless of power. Low power damages only the negative
result.
}

\formaldef{
The \textbf{power} of a test is
\[
  1 - \beta = \Prob(\text{reject } H_0 \mid H_1 \text{ true}) .
\]
It rises with the sample size $n$, with the size of the true effect, and with the
significance level $\alpha$ (at the cost of more false positives). It falls as the
data becomes noisier.

Conventionally $80\%$ power is treated as adequate. Below roughly $50\%$, a
non-rejection should not be reported as support for the null.
}

\analogybreak{
Power is computed against a \emph{specific} alternative. ``The test has 56\%
power'' is meaningless until the alternative is named --- here, a true exceedance
rate of $7\%$ against a claimed $5\%$. Against a true rate of $6\%$ the power is
much lower; against $15\%$, much higher. Any power figure is an answer to a
particular question.
}

\whyhere{
This is the reason downstream utility is evaluated on the pooled test set of about
$2{,}480$ observations rather than per market. It is also why QLIKE is computed
pooled: at per-market sample sizes the metric inverts, scoring Gaussian noise
above real data. Reporting a per-market Kupiec non-rejection as evidence that a
generator produces well-calibrated VaR would be a claim the data cannot support at
$17\%$ power.
}

%% ─────────────────────────────────────────────────────────────────────────────
\section{Multiple Testing}
\label{sec:multipletesting}

\situation{
Twenty researchers each test a different useless drug. Each uses the $5\%$
threshold. On average one of them finds a significant result and publishes it. The
drug does nothing; the statistics were applied correctly; the finding is still
spurious.
}

\problem{
This thesis computes many metrics per model per market. Every one is an
opportunity for a chance result, and the probability that \emph{at least one}
misleads grows quickly with their number.
}

\workedarith{
Each test has a $5\%$ false-positive rate under the null, so a $95\%$ chance of
behaving. For $m$ independent tests, the probability that all behave is $0.95^m$:

\medskip
\begin{tabular}{@{}lrr@{}}
\toprule
Tests $m$ & $\Prob(\text{no false positive})$ & $\Prob(\text{at least one})$ \\
\midrule
$1$  & $0.950$ & $0.050$ \\
$5$  & $0.774$ & $0.226$ \\
$10$ & $0.599$ & $0.401$ \\
$19$ & $0.377$ & $\mathbf{0.623}$ \\
\bottomrule
\end{tabular}

\medskip
At the 19 metrics of the original composite, the chance of at least one spurious
``significant'' result is $62\%$ --- more likely than not, with every individual
test correctly calibrated.

\medskip
\textbf{Bonferroni.} To hold the overall rate at $5\%$, test each at
$\alpha/m = 0.05/19 = 0.00263$. Then
$0.99737^{19} = 0.951$, so the family-wise error rate is about $4.9\%$. \checkmark

\medskip
\textbf{The price.} A threshold of $0.00263$ instead of $0.05$ demands far
stronger evidence per test, which drives power down --- and \S\ref{sec:power}
already showed power was the binding constraint. Correcting for multiplicity and
retaining power both require more data; there is no free version.
}

\formaldef{
With $m$ tests each at level $\alpha$, the \textbf{family-wise error rate} is
\[
  \mathrm{FWER} = 1 - (1-\alpha)^m
\]
under independence. The \textbf{Bonferroni correction} tests each at $\alpha/m$,
guaranteeing $\mathrm{FWER} \leq \alpha$ whether or not the tests are independent.

Less conservative alternatives --- Holm's step-down procedure, or
Benjamini--Hochberg control of the false discovery rate --- trade a weaker
guarantee for more power.
}

\analogybreak{
Bonferroni assumes the worst case and is conservative when tests are correlated,
which these emphatically are: \texttt{acf\_absolute\_mae} and
\texttt{acf\_squared\_mae} measure nearly the same thing, and Wasserstein and
energy distance both respond to the same distributional gaps. The effective number
of independent tests is well below 19, so a strict Bonferroni over-corrects.
Knowing the correction is too strong does not, however, license skipping it.
}

\whyhere{
This is part of why the evaluation framework ranks models rather than accumulating
significance tests. Ranks avoid the multiplicity problem: they make no
distributional claim per metric and cannot generate false positives to correct
for. The composite score is the mean of a fidelity rank and a temporal rank, and
the taxonomy in Chapter~\ref{ch:evaluation} exists precisely so that seven
permutation-invariant metrics cannot outvote seven ordering-sensitive ones ---
which is what let the shuffled control win the old unweighted 19-metric average
at $1.24$ against $2.47$.
}

%% ─────────────────────────────────────────────────────────────────────────────
\section{Linear Regression}
\label{sec:regression}

\situation{
Plot exam score against hours studied. The points scatter but trend upward. You
draw the straight line that best follows them, and use it to guess the score for
someone who studied six hours.
}

\problem{
``Best'' needs defining. Infinitely many lines pass near the points, and without a
criterion the choice is arbitrary --- and unrepeatable by anyone else.
}

\workedarith{
Four observations, hours $x$ and score $y$:

\medskip
\begin{tabular}{@{}lrrrr@{}}
\toprule
$x$ & 1 & 2 & 3 & 4 \\
$y$ & 2 & 4 & 5 & 9 \\
\bottomrule
\end{tabular}

\medskip
These are the numbers from \S\ref{sec:covariance}, where we found
$\bar{x} = 2.5$, $\bar{y} = 5$, $\Cov = 2.75$ and $\Var(x) = 1.25$.

\medskip
\textbf{Slope:}
\[
  \hat\beta = \frac{\Cov(x,y)}{\Var(x)} = \frac{2.75}{1.25} = 2.2 .
\]

\textbf{Intercept:} the line passes through $(\bar{x}, \bar{y})$, so
\[
  \hat\alpha = \bar{y} - \hat\beta\bar{x} = 5 - 2.2 \times 2.5 = -0.5 .
\]

Fitted line: $\hat{y} = -0.5 + 2.2x$.

\medskip
\textbf{Fitted values and residuals:}

\medskip
\begin{tabular}{@{}lrrrr@{}}
\toprule
$x$        & $1$   & $2$   & $3$   & $4$ \\
$y$        & $2$   & $4$   & $5$   & $9$ \\
$\hat{y}$  & $1.7$ & $3.9$ & $6.1$ & $8.3$ \\
residual   & $0.3$ & $0.1$ & $-1.1$ & $0.7$ \\
\bottomrule
\end{tabular}

\medskip
Residuals sum to $0.3 + 0.1 - 1.1 + 0.7 = 0$, as they must when an intercept is
fitted. \checkmark

Sum of squared residuals: $0.09 + 0.01 + 1.21 + 0.49 = 1.80$.

\medskip
\textbf{$R^2$.} Total squared deviation of $y$ about its mean:
$9 + 1 + 0 + 16 = 26$. So
\[
  R^2 = 1 - \frac{1.80}{26} = 0.931 .
\]
The line accounts for $93\%$ of the variation in scores.
}

\formaldef{
\textbf{Simple linear regression} models
\[
  y_i = \alpha + \beta x_i + \varepsilon_i ,
\]
and \textbf{ordinary least squares} chooses $\alpha, \beta$ to minimise
$\sum_i (y_i - \alpha - \beta x_i)^2$. Setting both partial derivatives
(\S\ref{sec:partials}) to zero gives
\[
  \hat\beta = \frac{\Cov(x,y)}{\Var(x)} , \qquad
  \hat\alpha = \bar{y} - \hat\beta \bar{x} .
\]
The \textbf{coefficient of determination} is
\[
  R^2 = 1 - \frac{\sum_i (y_i - \hat{y}_i)^2}{\sum_i (y_i - \bar{y})^2} ,
\]
the fraction of variance explained. In matrix form with several predictors,
$\hat{\boldsymbol\beta} = (X^{\!\top}X)^{-1}X^{\!\top}\mathbf{y}$ --- which exists
exactly when $X^{\!\top}X$ is invertible (\S\ref{sec:inverse}).
}

\analogybreak{
$R^2$ never decreases when predictors are added, even meaningless ones, so it
cannot be used to choose between models of different size. It also says nothing
about whether a linear model was appropriate: data lying on a perfect parabola can
produce a respectable $R^2$ from a straight line that is wrong everywhere. And
regression measures association, not causation --- the ice cream and drowning of
\S\ref{sec:covariance} would fit beautifully.
}

\whyhere{
$R^2$ is the diagnostic reported for the TimeGAN autoencoder: reconstruction $R^2$
above $0.8$ is required before Phase 3 is allowed to proceed, because a lossy
embedding means the adversarial stage trains on noise. The smoke run reached
$+0.995$. Regression also underlies the ARCH test, which regresses squared
residuals on their own lags, and the discriminative classifier, whose logistic
form is regression with a non-linear link.
}
